% Options for packages loaded elsewhere
\PassOptionsToPackage{unicode}{hyperref}
\PassOptionsToPackage{hyphens}{url}
%
\documentclass[
]{article}
\usepackage{amsmath,amssymb}
\usepackage{iftex}
\ifPDFTeX
  \usepackage[T1]{fontenc}
  \usepackage[utf8]{inputenc}
  \usepackage{textcomp} % provide euro and other symbols
\else % if luatex or xetex
  \usepackage{unicode-math} % this also loads fontspec
  \defaultfontfeatures{Scale=MatchLowercase}
  \defaultfontfeatures[\rmfamily]{Ligatures=TeX,Scale=1}
\fi
\usepackage{lmodern}
\ifPDFTeX\else
  % xetex/luatex font selection
\fi
% Use upquote if available, for straight quotes in verbatim environments
\IfFileExists{upquote.sty}{\usepackage{upquote}}{}
\IfFileExists{microtype.sty}{% use microtype if available
  \usepackage[]{microtype}
  \UseMicrotypeSet[protrusion]{basicmath} % disable protrusion for tt fonts
}{}
\makeatletter
\@ifundefined{KOMAClassName}{% if non-KOMA class
  \IfFileExists{parskip.sty}{%
    \usepackage{parskip}
  }{% else
    \setlength{\parindent}{0pt}
    \setlength{\parskip}{6pt plus 2pt minus 1pt}}
}{% if KOMA class
  \KOMAoptions{parskip=half}}
\makeatother
\usepackage{xcolor}
\usepackage{longtable,booktabs,array}
\usepackage{calc} % for calculating minipage widths
% Correct order of tables after \paragraph or \subparagraph
\usepackage{etoolbox}
\makeatletter
\patchcmd\longtable{\par}{\if@noskipsec\mbox{}\fi\par}{}{}
\makeatother
% Allow footnotes in longtable head/foot
\IfFileExists{footnotehyper.sty}{\usepackage{footnotehyper}}{\usepackage{footnote}}
\makesavenoteenv{longtable}
\setlength{\emergencystretch}{3em} % prevent overfull lines
\providecommand{\tightlist}{%
  \setlength{\itemsep}{0pt}\setlength{\parskip}{0pt}}
\setcounter{secnumdepth}{-\maxdimen} % remove section numbering
\ifLuaTeX
  \usepackage{selnolig}  % disable illegal ligatures
\fi
\usepackage{bookmark}
\IfFileExists{xurl.sty}{\usepackage{xurl}}{} % add URL line breaks if available
\urlstyle{same}
\hypersetup{
  pdftitle={ITSC 3155 Software Engineering Course Documentation},
  pdfauthor={Course Documentation},
  hidelinks,
  pdfcreator={LaTeX via pandoc}}

\title{ITSC 3155 Software Engineering Course Documentation}
\author{Course Documentation}
\date{}

\begin{document}
\maketitle

\textbf{Time and location}:~ Class is Asynchronous~Online and Office
Hours are via Zoom

\begin{center}\rule{0.5\linewidth}{0.5pt}\end{center}

\textbf{Instructor}:

\begin{itemize}
\tightlist
\item
  Dr. Mohsen Dorodchi
  \href{mailto:(Mohsen.Dorodchi@charlotte.edu)}{(\href{mailto:Mohsen.Dorodchi@charlotte.edu}{\nolinkurl{Mohsen.Dorodchi@charlotte.edu}})}

  \begin{itemize}
  \tightlist
  \item
    \textbf{Office hours}:

    \begin{itemize}
    \tightlist
    \item
      \emph{Woodward 430E or Woodward 432:}~Mondays 1:00 pm-2:00 pm or
      online (link)
    \item
      \emph{Online via zoom link in Canvas:}~Thursday 5:00pm-6:00pm
      (\href{https://charlotte-edu.zoom.us/j/98016945403}{Zoom LinkLinks
      to an external site.})~*~*
    \item
      If none of those times work, you can request for an appointment
      via email. Make sure to include the course number in the Subject
      Line~~and explain your question in the body of the email.
    \end{itemize}
  \end{itemize}
\end{itemize}

\begin{center}\rule{0.5\linewidth}{0.5pt}\end{center}

\textbf{Teaching Assistants}:

\begin{itemize}
\item
  Madison Melton
  \href{mailto:(mmelto21@charlotte.edu)}{(\href{mailto:mmelto21@charlotte.edu}{\nolinkurl{mmelto21@charlotte.edu}})}

  \begin{itemize}
  \tightlist
  \item
    Office Hours (Online)
    -~\href{https://charlotte-edu.zoom.us/j/92551802491?pwd=ZpMofRkZ2FuLt6LGIA5D0oeXaYX7cS.1}{Zoom
    linkLinks to an external site.}
  \item
    \textbf{Any changes in office hours or canceled sessions will be
    noted below during the applicable week}

    \begin{itemize}
    \tightlist
    \item
      Monday 3 PM to 5 PM
    \item
      Wednesday 3 PM to 5 PM~
    \item
      Friday 3 PM to 5 PM
    \end{itemize}
  \item
    \href{https://github.com/mmelto21-prog}{GitHub LinkLinks to an
    external site.}
  \end{itemize}
\item
  ~Geonwoo Lee (glee57@charlotte.edu)

  \begin{itemize}
  \item
    Office Hours~(Online)
    -~\href{https://charlotte-edu.zoom.us/j/9317544677}{Zoom linkLinks
    to an external site.}
  \item
    \textbf{Any changes in office hours or canceled sessions will be
    noted below during the applicable week}

    \begin{itemize}
    \tightlist
    \item
      Tuesday 10:00 AM \textasciitilde{} 12:00 PM (In-person office
      hours:~\textbf{Woodward Hall, Room 432}) \& Zoom available.
    \item
      Thursday 12:00 PM \textasciitilde{} 2:00 PM (Online)
    \item
      Friday 10:00 AM \textasciitilde{} 12:00 PM (Online)
      -\textgreater{} (**August 28, 1:00 PM to 3:00 PM)~**(Online)
    \end{itemize}
  \item
    \href{https://github.com/Jamesonthehill}{GitHub LinkLinks to an
    external site.}
  \end{itemize}
\end{itemize}

\begin{center}\rule{0.5\linewidth}{0.5pt}\end{center}

\href{https://instructure.charlotte.edu/courses/264138/files/33808309/preview}{Course
Overview}

\begin{center}\rule{0.5\linewidth}{0.5pt}\end{center}

Online books:

\begin{itemize}
\tightlist
\item
  Winters, Titus, Tom Manshreck, and Hyrum
  Wright.~\href{https://abseil.io/resources/swe-book}{Software
  Engineering at GoogleLinks to an external site.}: Lessons learned from
  programming over time. O\textquotesingle Reilly Media, 2020.
\item
  O\textquotesingle Regan,
  Gerard.~\href{https://link.springer.com/content/pdf/10.1007/978-3-319-57750-0.pdf}{Concise
  guide to software engineeringLinks to an external site.}. Springer,
  Cham, 2017.
\end{itemize}

\begin{center}\rule{0.5\linewidth}{0.5pt}\end{center}

Click on a module button to view what\textquotesingle s happening in
that week.

\begin{longtable}[]{@{}cccll@{}}
\toprule\noalign{}
\textbf{Module} & \textbf{Topic} & \textbf{Homework} & & \\
\midrule\noalign{}
\endhead
\bottomrule\noalign{}
\endlastfoot
\textbf{Introduction} & \textbf{Course Overview} & ~ & & \\
\textbf{Week 1} & \textbf{Introduction to Software Engineering} &
\href{https://instructure.charlotte.edu/courses/264138/modules/items/8170989}{\textbf{Python
Crash Course}} &
\href{https://instructure.charlotte.edu/courses/264138/assignments/2981563}{\textbf{Environment
Setup}} & \\
\textbf{Intro to Python} & & & & \\
\textbf{Week 2} & \textbf{Requirement Engineering} &
\href{https://instructure.charlotte.edu/courses/264138/pages/python-crash-course-part2?wrap=1}{Python
Crash Course 2} &
\href{https://instructure.charlotte.edu/courses/264138/assignments/2981587?wrap=1}{Assignment
0 Python Class} &
\href{https://instructure.charlotte.edu/courses/264138/quizzes/636082?wrap=1}{Quiz
- Requirements} \\
\textbf{OOP in Practice - 1} & & & & \\
\textbf{Week 3} & \textbf{Software Design} &
\href{https://instructure.charlotte.edu/courses/264138/assignments/2981548?wrap=1}{Assignment1
- Sandwich Maker} &
\href{https://instructure.charlotte.edu/courses/264138/quizzes/636064?wrap=1}{Quiz
- Software Design} & \\
\textbf{Version Control} & & & & \\
\textbf{Week 4} & \textbf{Agile Software Development} &
\href{https://instructure.charlotte.edu/courses/264138/assignments/2981550?wrap=1}{Assignment
2 - Modular Sandwich Maker} &
\href{https://instructure.charlotte.edu/courses/264138/assignments/2981512?wrap=1}{Quiz
- Agile Software Development} & \\
\textbf{OOP in Practice - 2} & & & & \\
\textbf{Week 5} & **Software Implementation and Deployment~** &
\href{https://instructure.charlotte.edu/courses/264138/assignments/2981552?wrap=1}{Assignment
3 - Sandwich Maker Database} &
\href{https://instructure.charlotte.edu/courses/264138/assignments/2981534?wrap=1}{Quiz
- Software Implementation} & \\
\textbf{Intro to Relational Database} & & & & \\
\textbf{Week 6} & **Application Programming Interface (API)~** &
\href{https://instructure.charlotte.edu/courses/264138/assignments/2981553?wrap=1}{Assignment
4 - FastAPI Implementation} &
\href{https://instructure.charlotte.edu/courses/264138/assignments/2981502?wrap=1}{Quiz
- Software Testing} & \\
**Software Testing~** & & & & \\
\textbf{Midterm} & \textbf{Midterm-Prep} & Midterm & & \\
\textbf{Week 7} & \textbf{Scrum In Depth} &
\href{https://instructure.charlotte.edu/courses/264138/assignments/2981554?wrap=1}{Assignment
5 - CRUD Operation in FastAPI} &
\href{https://instructure.charlotte.edu/courses/264138/quizzes/636022?wrap=1}{Quiz
- Security} & \\
\textbf{TDD / FastAPI- Part2} & & & & \\
\textbf{Week 8} & \textbf{CRUD Operations Using FastAPI} &
\href{https://instructure.charlotte.edu/courses/264138/quizzes/636043?wrap=1}{Quiz
- Scrum} & & \\
\textbf{Security} & & & & \\
\textbf{Week 9} & \textbf{Frontend - Web Development} &
\href{https://instructure.charlotte.edu/courses/$CANVAS_OBJECT_REFERENCE$/assignments/g53104bb09d2c7e6e864e6cfc9c4d538b?wrap=1}{Group
Project - Part 1} & & \\
\textbf{Week 10} & \textbf{Containerization} & \textbf{Docker} &
\href{https://instructure.charlotte.edu/courses/264138/quizzes/636027?wrap=1}{Quiz
- Containerization {[}OPTIONAL{]}} & \\
\textbf{Week 11} & \textbf{DevOps} &
\href{https://instructure.charlotte.edu/courses/$CANVAS_OBJECT_REFERENCE$/assignments/ge95cf156bf9319553487d7b054086722?wrap=1}{Group
Project - Part 2} &
\href{https://instructure.charlotte.edu/courses/264138/quizzes/636032?wrap=1}{Quiz
- DevOps {[}Optional{]}} & \\
\textbf{Group Project Meeting} & \textbf{GitHub Actions} & & & \\
\textbf{Week 12} & \textbf{Cloud Computing} &
\href{https://instructure.charlotte.edu/courses/264138/quizzes/636063?wrap=1}{Quiz
- Cloud Computing {[}Optional{]}} & ~ & \\
\textbf{Group Project Meeting} & \textbf{SQLAlchemy Join} & & & \\
\textbf{Week 13} & **AI/Machine Learning I~** & Group Project & & \\
\textbf{Week 14/15/16} & \textbf{AI/Machine Learning II} & Group Project
& & \\
\end{longtable}

\begin{center}\rule{0.5\linewidth}{0.5pt}\end{center}

\textbf{Grading and Assessments}

Grading will be based on the following scale:

\begin{itemize}
\tightlist
\item
  A~ 90\% - 100\%
\item
  B~ 80\% - 89\%
\item
  C~ 70\% - 79\%
\item
  D~ 60\% - 69\%
\item
  F~ 0\% - 59\%
\end{itemize}

You can find the breakdown
in~\href{https://instructure.charlotte.edu/courses/$CANVAS_COURSE_REFERENCE$/modules/items/g12bc07101b82cdbb0110e8daf740ba16}{Course
Overview~}. More details about each of the assessment activities is
available on the Canvas course site.

\begin{center}\rule{0.5\linewidth}{0.5pt}\end{center}

\textbf{Academic integrity policy}:

\begin{itemize}
\tightlist
\item
  \textbf{All work must be the student's own}. All external references
  used in reports must be properly cited. No credit will be given for
  duplicate or plagiarized work.
\item
  Faculty may ask students to produce identification at examinations and
  may require students to demonstrate that graded assignments completed
  outside of class are their own work.
\item
  All students are required to read and abide by
  the~\href{https://legal.uncc.edu/policies/up-407}{Code of Student
  Academic Integrity}
\end{itemize}

Violations of the Code of Student Academic Integrity, including
plagiarism, will result in disciplinary action as provided in the Code.
Definitions and examples of plagiarism are set forth in the Code and on
the~\href{https://scai.uncc.edu/academic-integrity/academic-misconduct-policies-examples-0}{Student
Conduct and Academic Integrity website~}.

\begin{center}\rule{0.5\linewidth}{0.5pt}\end{center}

\textbf{Tutoring Center}

CCI Undergraduate Tutoring Center offers ~free, one-on-one tutoring
services to all undergraduate students who request their service. You
can make an appointment and also view tutors\textquotesingle{}
availability for drop-ins via Connect. Here are instructions
on~\href{https://cci.charlotte.edu/student-life/tutoring}{how you can do
this}.~

\begin{center}\rule{0.5\linewidth}{0.5pt}\end{center}

\textbf{Other policies}:

\begin{itemize}
\tightlist
\item
  \textbf{Mutual respect and non-discrimination}: All students and the
  instructor are expected to engage with each other respectfully.
  Unwelcome conduct directed toward another person may constitute a
  violation of~\href{https://legal.uncc.edu/policies/up-406}{The Code of
  Student Responsibility}
\item
  Any student suspected of engaging in such conduct will be referred to
  the~\href{https://scai.uncc.edu/}{Office of Student Conduct.}
\item
  \textbf{Disability-based accommodations}: Students in this course
  seeking disability-based accommodations must first consult with the
  Office of Disability Services and follow the instructions of that
  office for obtaining accommodations.
\item
  \textbf{Mental health concerns}~or stressful events may reduce a
  student\textquotesingle s ability to participate in daily activities
  or diminish academic performance. Services are available to assist you
  with addressing these and other concerns you may be experiencing. You
  can learn more about the broad range of confidential mental health
  services available on campus via the Counseling \& Psychological
  Services (CAPS) website
  at~\href{http://caps.uncc.edu/}{caps.charlotte.edu}
\item
  \textbf{Extenuating circumstances}~for absences and missed exams
  include medical appointments, military/court orders, and/or personal
  and family emergencies, such as a death in the immediate family, where
  a student is able to provide an instructor with appropriate supporting
  documentation of the absence.

  \begin{itemize}
  \tightlist
  \item
    The Office of Student Assistance and Support Services (SASS) may
    also be able to assist with verification of such emergencies, once a
    student is able to return to classes.
  \item
    In cases of absence due to pregnancy or parenting, students should
    contact the Title IX Office to obtain absence verification by
    completing the form
    at~{[}\url{http://bit.ly/332eaGd\%5D(http://bit.ly/332eaGd)}
  \end{itemize}
\end{itemize}

\begin{center}\rule{0.5\linewidth}{0.5pt}\end{center}

\section{Assignment 1 - Sandwich
Maker}\label{assignment-1---sandwich-maker}

This assignment has two parts.

\textbf{Part1 (40 points)}

As part of this assignment, we will create a Ham Sandwich Maker Machine
using interactive Python terminal. An automatic ham sandwich maker uses
a set of programmed ingredients to make sandwiches. First, users select
the size of the sandwich, then the machine determines how much resources
are available before making it. The user receives a message if there are
insufficient resources. Otherwise, the user must pay for sandwich by
coins. The machine will indicate if insufficient coins are inserted or
if there is an exchange amount when coins are inserted. Also, on the
screen, users can view the remaining resources report and turn off the
machine.

\begin{enumerate}
\def\labelenumi{\arabic{enumi}.}
\tightlist
\item
  Read and analyze the project document carefully:~
\end{enumerate}

\begin{itemize}
\tightlist
\item
\end{itemize}

\begin{itemize}
\tightlist
\item
  \href{https://instructure.charlotte.edu/courses/264138/files/33808382?wrap=1}{Sandwich
  Maker Machine-
  Document.pdf}\href{https://instructure.charlotte.edu/courses/264138/files/33808382/download?download_frd=1}{Download
  Sandwich Maker Machine- Document.pdf}
\end{itemize}

2. Download the code skeleton:

\begin{itemize}
\tightlist
\item
\end{itemize}

\begin{itemize}
\tightlist
\item
  \href{https://github.com/samfallahian/ITSC3155-Assignments}{Source
  codeLinks to an external site.}
\end{itemize}

3. Open the project in PyCharm and complete the missing parts.

4. You must submit the~\textbf{main.py}~file for this part.

\textbf{Part2 (10 points)}

You should apply your knowledge of Git to this project and push your
project into your GitHub account by following the instructions below:

\begin{itemize}
\tightlist
\item
  \href{https://docs.google.com/document/d/1DahT-ZiFTmzyIGTBtEwgMe0WFDMs5ydjbg0b-U2yRfs/edit}{\textbf{Sandwich
  Maker Machine-Version ControlLinks to an external site.}}
\end{itemize}

Submit a~**link~**to your private GitHub project.

\textbf{Turn in:}

\begin{itemize}
\tightlist
\item
  main.py file
\item
  GitHub link
\end{itemize}

\textbf{Rubrics:}

\textbf{Part 1:}

\begin{longtable}[]{@{}cc@{}}
\toprule\noalign{}
Criteria & Points \\
\midrule\noalign{}
\endhead
\bottomrule\noalign{}
\endlastfoot
Program runs properly based on provided scenario (Sample interactive
terminal input/output) without hard-coding. & 8 \\
\end{longtable}

\begin{longtable}[]{@{}cccccc@{}}
\toprule\noalign{}
Function & Input & Points & Output & Points & Works? \\
\midrule\noalign{}
\endhead
\bottomrule\noalign{}
\endlastfoot
check\_resources() & ingredients (dictionary) & 2 & True/False (boolean)
& 2 & 4 \\
process\_coins() & NAN & 2 & Number (float) & 2 & 4 \\
transaction\_result() & coins (float), cost (float) & 2 & True/False
(boolean) & 2 & 4 \\
make\_sandwich() & sandwich\_size (string), ingredients (dictionary) & 2
& NAN & 2 & 4 \\
\end{longtable}

Part 2:

\begin{longtable}[]{@{}cc@{}}
\toprule\noalign{}
Criteria & Points \\
\midrule\noalign{}
\endhead
\bottomrule\noalign{}
\endlastfoot
The source code has been uploaded to your GitHub account. & 2 \\
The code is pushed into the repository (manual uploading is not
acceptable). & 3 \\
The GitHub repo history shows at least five commits. & 3 \\
Commit messages follow the conventional "conventional commits" format. &
2 \\
\end{longtable}

Points

50

Submitting

a website url or a file upload

File Types

py

\begin{longtable}[]{@{}llll@{}}
\toprule\noalign{}
DueForAvailable fromUntil & & & \\
\midrule\noalign{}
\endhead
\bottomrule\noalign{}
\endlastfoot
Sep 13Sep 13 at 11:59pm & Everyone & -N/A & -N/A \\
\end{longtable}

\begin{center}\rule{0.5\linewidth}{0.5pt}\end{center}

\section{Assignment 2 - Modular Sandwich
Maker}\label{assignment-2---modular-sandwich-maker}

In this assignment, we will convert the sandwich maker code (assignment
1) into a modular format. In the following file, you will find
instructions on how to proceed:

\href{https://instructure.charlotte.edu/courses/264138/files/33808302?wrap=1}{Modular
Sandwich Maker
Machine.pdf}\href{https://instructure.charlotte.edu/courses/264138/files/33808302/download?download_frd=1}{Download
Modular Sandwich Maker Machine.pdf}

\href{https://github.com/samfallahian/ITSC3155-Assignments}{\emph{Download
code skeleton here.Links to an external site.}}

Turn in:

\textbf{Your GitHub Repository link} \textbf{Make sure you commit and
push the files to GitHub (DO NOT SUBMIT A ZIP FILE)}

Points

30

Submitting

a website url

\begin{longtable}[]{@{}llll@{}}
\toprule\noalign{}
DueForAvailable fromUntil & & & \\
\midrule\noalign{}
\endhead
\bottomrule\noalign{}
\endlastfoot
Sep 27Sep 27 at 11:59pm & Everyone & & \\
\end{longtable}

Modular Ham Sandwich Maker Machine Program In this assignment, you will
use what you have learned about Python modules to make our Ham Sandwich
Maker Machine code more object-oriented. The program should be split
into four separate files. To complete this assignment, follow these
steps. 1- Code Skeleton: • On Canvas, you can find the code skeleton on
the assignment page. 2- Import all modules into main.py • Three modules
(data, sandwich\_maker, cashier) should be imported at the top of
main.py. ``import \textless file\_name\textgreater'' • Create two
variables based on data dictionaries (resources and recipes). • Then you
need to create instance from each. Pay attention sandwich\_maker class
has a constructor variable (resources) which you imported in last step.
``\textless var\textgreater{} =
\textless imported\_name\textgreater.\textless class\_name\textgreater''
3- Place functions in corresponding module: • The simplest way to
accomplish this is to copy and paste the function into the corresponding
module according to the instructions at the end of this document. 4-
Fine tune main.py • Change the variable names if necessary and run the
program. To test your program, you can use the same scenario as in
assignment 1. Project Structure: SandwichMaker class in
sandwich\_makerr.py file: • This class contains everything about making
sandwiches. The check\_resources() and make\_sandwich() functions should
be placed in this class. Cashier class in cashier.py file: • This class
contains everything about purchasing. The process\_coins() and
transaction\_result() functions should be placed in this class. Data.py:
• This file keeps the data we need to run program like a database! Later
in the course, we will convert it into a Python class that stores and
retrieves data. Main.py: • It acts as the starting point of execution
for the program.

\begin{center}\rule{0.5\linewidth}{0.5pt}\end{center}

\section{Assignment 3 - Sandwich Maker
Database}\label{assignment-3---sandwich-maker-database}

For this assignment you need to create a database and table for Sandwich
Maker Machine project and insert some dummy data. Follow the instruction
below:

1- Install MySQL and MySQL Workbench:
(\href{https://www.youtube.com/watch?v=7S_tz1z_5bA}{MySQL Tutorial for
Beginners by Mosh HamedaniLinks to an external site.})

2- Follow
the~\href{https://instructure.charlotte.edu/courses/264138/files/33808460?wrap=1}{instruction~}\href{https://instructure.charlotte.edu/courses/264138/files/33808460/download?download_frd=1}{Download
instruction}and create 3 tables and insert all data that you used in
Assignment 1 into corresponding tables.

3- Turn in:

\begin{itemize}
\tightlist
\item
  \textbf{Three screenshots~(result of SELECT statement) for each table
  based on the instructions.}
\item
  \textbf{A .sql file containing all queries}
\end{itemize}

Points

50

Submitting

a text entry box, a website url, or a file upload

\begin{longtable}[]{@{}llll@{}}
\toprule\noalign{}
DueForAvailable fromUntil & & & \\
\midrule\noalign{}
\endhead
\bottomrule\noalign{}
\endlastfoot
Oct 5Oct 5 at 11:59pm & Everyone & -N/A & -N/A \\
\end{longtable}

~Assignment 3 - Sandwich Maker Database Rubric

\begin{center}\rule{0.5\linewidth}{0.5pt}\end{center}

Enable self assessment

\begin{center}\rule{0.5\linewidth}{0.5pt}\end{center}

\section{Assignment 4 - FastAPI
Implementation}\label{assignment-4---fastapi-implementation}

The goal of this assignment is to create a simple API using FastAPI
framework. Using the instruction below for doing your assignment.~**Make
sure to regularly commit to your remote repository. At least 3 are
required.~**

\textbf{!!!!MAKE SURE TO INVITE AS COLLABORATORS THE TAs TO THE GITHUB
REPOSITORY!!!!}

\href{https://docs.google.com/document/d/1JaDI09XzxiohC0ko6v3VS9BAPWd8l-pOrR7Y49KN1eM/edit?tab=t.0}{FastAPI-
ImplementationLinks to an external site.}

Turn in:

\begin{enumerate}
\def\labelenumi{\arabic{enumi}.}
\tightlist
\item
  \textbf{Your GitHub repository that shows your work (complete source
  code).}

  \begin{itemize}
  \tightlist
  \item
    \textbf{py file should be included in assignment folder.}
  \item
    \textbf{Your commit messages must follow conventional commits.}
  \item
    \textbf{Your last commit time must be before your submission date}
  \end{itemize}
\item
  \textbf{Screenshots of the results taken from the Interactive API docs
  (Part 1-step 5, Part2-step 2).}

  \begin{itemize}
  \tightlist
  \item
    \textbf{Run~ALL~requests in each part.}
  \end{itemize}
\end{enumerate}

Points

15

Submitting

a text entry box, a website url, or a file upload

\begin{longtable}[]{@{}llll@{}}
\toprule\noalign{}
DueForAvailable fromUntil & & & \\
\midrule\noalign{}
\endhead
\bottomrule\noalign{}
\endlastfoot
Oct 23Oct 23 at 11:59pm & Everyone & -N/A & -N/A \\
\end{longtable}

\begin{center}\rule{0.5\linewidth}{0.5pt}\end{center}

\section{Assignment 5 - CRUD Operation in FastAPI - Sandwich Maker
Machine
Program}\label{assignment-5---crud-operation-in-fastapi---sandwich-maker-machine-program}

In this assignment, you will have the opportunity to explore and
implement the fundamental principles of creating, reading, updating, and
deleting (CRUD) operations on a database using FastAPI, a modern and
highly efficient web framework for building APIs with Python.

The primary goal of this assignment is to develop a comprehensive
understanding of CRUD operations by implementing them on a database.
Specifically, you will be working with an extended version of a Sandwich
Maker Machine Program. Throughout this assignment, you will become
familiar with the following key concepts and techniques:

\begin{itemize}
\tightlist
\item
  \textbf{Modular Implementation with FastAPI:}~You will learn how to
  structure your code in a modular fashion using FastAPI, ensuring that
  your application is organized, maintainable, and extensible.
\item
  \textbf{Database Connectivity:}~You will explore how to establish a
  connection to a database from within a FastAPI application. This
  connection will serve as the foundation for storing and retrieving
  data related to our Sandwich Maker Machine Program.
\item
  \textbf{Table Creation with FastAPI:}~You will dive into the process
  of defining database tables within your FastAPI application,
  specifying the schema, and handling migrations if necessary.
\end{itemize}

\textbf{Endpoint Creation:}~You will create API endpoints for performing
CRUD operations on the database. These endpoints will allow you to
interact with the Sandwich Maker Machine Program by adding, retrieving,
updating, and deleting sandwich-related data.

Read the assignment instruction~several~time
:~\href{https://instructure.charlotte.edu/courses/264138/files/33808321?wrap=1}{CRUD
Operation in FastAPI - Sandwich Maker Machine
Program}\href{https://instructure.charlotte.edu/courses/264138/files/33808321/download?download_frd=1}{Download
CRUD Operation in FastAPI - Sandwich Maker Machine Program}

\href{https://github.com/samfallahian/ITSC3155-Assignments}{Download
code here.Links to an external site.}

Turn in:

\begin{enumerate}
\def\labelenumi{\arabic{enumi}.}
\tightlist
\item
  \textbf{Your GitHub repository that shows your work (complete source
  code).}

  \begin{itemize}
  \tightlist
  \item
    \textbf{Files should be included in assignment folder.}
  \item
    \textbf{Your commit messages must follow conventional commits.}
  \item
    \textbf{Your last commit time must be before your submission date.}
  \end{itemize}
\end{enumerate}

Points

80

Submitting

a website url

\begin{longtable}[]{@{}llll@{}}
\toprule\noalign{}
DueForAvailable fromUntil & & & \\
\midrule\noalign{}
\endhead
\bottomrule\noalign{}
\endlastfoot
Oct 30Oct 30 at 11:59pm & Everyone & -N/A & \\
\end{longtable}

\begin{center}\rule{0.5\linewidth}{0.5pt}\end{center}

\section{Group Project - Part 1 - Requirement Analysis and User
Stories}\label{group-project---part-1---requirement-analysis-and-user-stories}

\textbf{GitHub Repository:}

\begin{itemize}
\tightlist
\item
  \textbf{Create a new clean repository (One of you).}
\item
  \textbf{Invite your group member as well as Instructor team (Professor
  and TAs)}
\item
  \textbf{Download
  the~\textbf{\href{https://github.com/samfallahian/ITSC3155-Assignments/tree/main/FinalProject}{\textbf{code
  skeletonLinks to an external site.}}}~and push it to the repository
  with a proper commit message}
\item
  **The rest of the group should clone the project
  using~~**\textbf{\texttt{git\ clone\ \textless{}repo\_address\textgreater{}}}
\end{itemize}

\textbf{Project Setup:}

\begin{itemize}
\tightlist
\item
  \textbf{Create a Virtual Environment and
  install~\textbf{\href{https://github.com/samfallahian/ITSC3155-Assignments/blob/main/FinalProject/README.md}{\textbf{dependenciesLinks
  to an external site.}}}~.}
\item
  \textbf{Run the FastApi server and test the sample code.}
\end{itemize}

\textbf{Requirement Analysis and User Stories:}

\begin{itemize}
\tightlist
\item
  \textbf{Pay close attention to what the client wants
  in~\textbf{\href{https://instructure.charlotte.edu/courses/264138/files/33808470?wrap=1}{\textbf{this
  document}}\href{https://instructure.charlotte.edu/courses/264138/files/33808470/download?download_frd=1}{\textbf{Download
  this document}}}.}
\item
  **Extract requirements and fill
  out~**\href{https://instructure.charlotte.edu/courses/264138/files/33808457?wrap=1}{\textbf{Product
  Backlogs.}}
\item
  \textbf{Create~\textbf{\href{https://instructure.charlotte.edu/courses/264138/files/33808471?wrap=1}{\textbf{User
  Stories}}\href{https://instructure.charlotte.edu/courses/264138/files/33808471/download?download_frd=1}{\textbf{Download
  User Stories}}}~for your backlogs.}
\end{itemize}

\textbf{Deliverables:}

\begin{itemize}
\tightlist
\item
  \textbf{GitHub repository link (With code skeleton and all
  collaborators).}
\item
  \textbf{A document
  for~\textbf{\href{https://instructure.charlotte.edu/courses/264138/files/33808457?wrap=1}{\textbf{Product
  Backlogs}}}.}
\item
  \textbf{A document
  for~\textbf{\href{https://instructure.charlotte.edu/courses/264138/files/33808471?wrap=1}{\textbf{User
  Stories}}}.}
\end{itemize}

\textbf{Grading Criteria:}

\begin{itemize}
\tightlist
\item
  **~~The GitHub repository has been set up correctly (collaborators,
  code skeleton).**
\item
  **~Conventional commit messages.**
\item
  \textbf{Quality of product backlog and user stories that meets all
  customer needs.}
\item
  \textbf{Attendance (participation in group meetings).}
\end{itemize}

**~One submission per group is enough.**

Points

15

Submitting

a text entry box, a website url, or a file upload

\begin{longtable}[]{@{}llll@{}}
\toprule\noalign{}
DueForAvailable fromUntil & & & \\
\midrule\noalign{}
\endhead
\bottomrule\noalign{}
\endlastfoot
Nov 2Nov 2 at 11:59pm & Everyone & -N/A & -N/A \\
\end{longtable}

\begin{center}\rule{0.5\linewidth}{0.5pt}\end{center}

\section{Group Project - Part 2- Software and Database
Design}\label{group-project---part-2--software-and-database-design}

\textbf{Instruction:}

\begin{enumerate}
\def\labelenumi{\arabic{enumi}.}
\item
  \textbf{Follow
  this~\textbf{\href{https://instructure.charlotte.edu/courses/264138/files/33808474?wrap=1}{**instruction~**}\href{https://instructure.charlotte.edu/courses/264138/files/33808474/download?download_frd=1}{\textbf{Download
  instruction}}}~and create at list~three~different block diagram to
  show your design. Your design diagrams won\textquotesingle t be graded
  against strict UML rules, creating casual diagrams is enough. However,
  they help you understand and communicate your design effectively.
  Focus on clarity and simplicity rather than perfection.}
\item
  **After deciding about your database tables, implement models and
  schemas. You should have separate file for each table inside models
  and schemas packages. Don\textquotesingle t forget to
  update~\textbf{\href{https://github.com/samfallahian/ITSC3155-Assignments/blob/main/FinalProject/api/models/model_loader.py}{\textbf{\texttt{model\_loader.py.}}}\href{https://github.com/samfallahian/ITSC3155-Assignments/blob/main/FinalProject/api/models/model_loader.py}{\textbf{Links
  to an external site.}}}~**

  \begin{itemize}
  \tightlist
  \item
    **Database Design Consideration:~**Read
    this~\href{https://instructure.charlotte.edu/courses/264138/files/33808129?wrap=1}{\textbf{document}}\href{https://instructure.charlotte.edu/courses/264138/files/33808129?wrap=1}{~}\href{https://instructure.charlotte.edu/courses/264138/files/33808129/download?download_frd=1}{Download
    document}carefully and make sure you covered the minimum
    requirements.
  \item
    **Primary Key Requirement:~**Ensure that every table in your
    database has a primary key.
  \item
    **Clear Table Relationships:~**It's important that the relationships
    between tables are clearly implemented and defined.
  \end{itemize}
\item
  **If necessary, update
  your~**\href{https://instructure.charlotte.edu/courses/264138/files/33808457?wrap=1}{\textbf{Product
  Backlogs}}\href{https://instructure.charlotte.edu/courses/264138/files/33808457/download}{**~**}**and~**\href{https://instructure.charlotte.edu/courses/264138/files/33808471?wrap=1}{\textbf{User
  Stories.}}
\end{enumerate}

\textbf{Code Skeleton:}

The code skeleton I provided contains sample tables and endpoints that
are not specifically tailored to our project\textquotesingle s
requirements. These examples were included to demonstrate how various
modules and components of an API can work together in a real-world
application. They are intended to serve as a reference or starting point
for understanding the structure and integration of different parts of an
API. Therefore, you are not required to use the exact table structures
or endpoints as they are in the code skeleton. I encourage you to design
and implement your database tables and API endpoints based on the
requirements you have extracted from the project brief. Your focus
should be on addressing the specific needs of the Online Restaurant
Ordering System.

\textbf{Deliverables:}

\begin{itemize}
\item
  \textbf{A pdf file containing at least three casual block diagrams.}
\item
  \textbf{GitHub repository link (With updated codes before assignment
  due date and conventional commit messages).}
\item
  \textbf{Updated document
  for~\textbf{\href{https://instructure.charlotte.edu/courses/264138/files/33808457?wrap=1}{\textbf{Product
  Backlogs}}}. (If any)}
\item
  \textbf{Updated document
  for~\textbf{\href{https://instructure.charlotte.edu/courses/264138/files/33808471?wrap=1}{\textbf{User
  Stories}}}. (If any)}
\end{itemize}

\textbf{Grading Criteria:}

\begin{itemize}
\tightlist
\item
  \textbf{Quality of software design (block diagrams).}
\item
  \textbf{Your code is correctly implemented like code skeleton.}
\item
  \textbf{All tables must be created after running the API.}
\item
  \textbf{Attendance (participation in group meetings).}
\end{itemize}

**~One submission per group is enough.**

Points

20

Submitting

a text entry box, a website url, or a file upload

\begin{longtable}[]{@{}llll@{}}
\toprule\noalign{}
DueForAvailable fromUntil & & & \\
\midrule\noalign{}
\endhead
\bottomrule\noalign{}
\endlastfoot
Nov 18Nov 18 at 11:59pm & Everyone & -N/A & -N/A \\
\end{longtable}

\begin{center}\rule{0.5\linewidth}{0.5pt}\end{center}

\section{Group Project Part 3 - Final
Product}\label{group-project-part-3---final-product}

\textbf{Based on your requirement analysis, you must work on~two
sprints~and deliver a completed software (API) to your customer by the
deadline specified in this assignment.}

\textbf{Instruction:}

\begin{enumerate}
\def\labelenumi{\arabic{enumi}.}
\item
  \textbf{Select user stories from your Product Backlogs and begin
  breaking them down into tasks for developers fill
  out~~\textbf{\href{https://instructure.charlotte.edu/courses/264138/files/33808466?wrap=1}{\textbf{Sprint
  Backlog -
  Tasks}}\href{https://instructure.charlotte.edu/courses/264138/files/33808466/download?download_frd=1}{\textbf{Download
  Sprint Backlog - Tasks}}}.}

  \begin{enumerate}
  \def\labelenumii{\arabic{enumii}.}
  \tightlist
  \item
    \textbf{Include the user stories you want to include in Sprint 1}

    \begin{enumerate}
    \def\labelenumiii{\arabic{enumiii}.}
    \tightlist
    \item
      We recommend your first sprint should~AT MINIMUM~cover the
      following timeframe:~April 15 - 22)
    \end{enumerate}
  \item
    \textbf{Include the user stories you want to include in Sprint 2}

    \begin{enumerate}
    \def\labelenumiii{\arabic{enumiii}.}
    \tightlist
    \item
      We recommend your first sprint should~AT MINIMUM~cover the
      following timeframe:~April 23 - May 1 )
    \end{enumerate}
  \end{enumerate}
\item
  \textbf{For each task, you need to implement corresponding endpoint
  (or endpoints) and make sure your implementation can be traced.}
\item
  \textbf{If during development, you discover that your product backlogs
  or user stories need to be updated, reordered, or added, update the
  corresponding files
  (}\href{https://instructure.charlotte.edu/courses/264138/files/33808471?wrap=1}{\textbf{User
  Stories}}\href{https://instructure.charlotte.edu/courses/264138/files/33808471/download?download_frd=1}{\textbf{Download
  User
  Stories}}**~and~\textbf{\href{https://instructure.charlotte.edu/courses/264138/files/33808457?wrap=1}{\textbf{Product
  Backlog}}\href{https://instructure.charlotte.edu/courses/264138/files/33808457/download?download_frd=1}{\textbf{Download
  Product Backlog}}}).**
\item
  \textbf{Implement~at least one unit test using pytest.}
\item
  \textbf{Make sure everything works properly by testing your endpoints
  and running unit tests.}
\item
  \textbf{At the end of each sprint, fill
  out~\textbf{\href{https://instructure.charlotte.edu/courses/264138/files/33808472?wrap=1}{\textbf{Sprint
  Review}}\href{https://instructure.charlotte.edu/courses/264138/files/33808472/download?download_frd=1}{\textbf{Download
  Sprint Review}}}~form for each sprint.}
\end{enumerate}

\textbf{Deliverables:}

\begin{itemize}
\item
  \textbf{A document for both sprint by filling out
  the~\textbf{\href{https://instructure.charlotte.edu/courses/264138/files/33808466?wrap=1}{\textbf{Sprint
  Backlog -
  Tasks}}\href{https://instructure.charlotte.edu/courses/264138/files/33808466/download?download_frd=1}{\textbf{Download
  Sprint Backlog - Tasks}}}.}

  \begin{itemize}
  \tightlist
  \item
    \textbf{Include sprint 1}
  \item
    \textbf{Include sprint 2}
  \end{itemize}
\item
  \textbf{GitHub repository link (contains all of the necessary files
  and code for running the project and unit test pushed before
  assignment due date and conventional commit messages).}
\item
  \textbf{Two documents for~sprint review
  1~\emph{\textbf{\textbf{AND}}}~2~by filling
  out~\textbf{\href{https://instructure.charlotte.edu/courses/264138/files/33808472?wrap=1}{\textbf{Sprint
  Review.docx}}\href{https://instructure.charlotte.edu/courses/264138/files/33808472/download?download_frd=1}{\textbf{Download
  Sprint Review.docx}}}.}

  \begin{itemize}
  \tightlist
  \item
    \textbf{Include sprint 1 review}
  \item
    \textbf{Include sprint 2 review}
  \end{itemize}
\item
  \textbf{Updated document
  for~\textbf{\href{https://instructure.charlotte.edu/courses/264138/files/33808457?wrap=1}{\textbf{Product
  Backlogs}}}. (If any)}
\item
  \textbf{Updated document
  for~\textbf{\href{https://instructure.charlotte.edu/courses/264138/files/33808471?wrap=1}{\textbf{User
  Stories}}}. (If any)}
\end{itemize}

\textbf{Evaluation Checklist:}

\begin{itemize}
\tightlist
\item
  \textbf{Your final product should have CRUD operation for all tables.}
\item
  \textbf{Your final product must at least
  address~\textbf{\href{https://instructure.charlotte.edu/courses/264138/files/33808478?wrap=1}{\textbf{these
  questions}}\href{https://instructure.charlotte.edu/courses/264138/files/33808478/download?download_frd=1}{\textbf{Download
  these questions}}}. This list can be used as a checklist to evaluate
  your final product.}
\item
  \textbf{Your final product must work properly based on instruction.}
\end{itemize}

\textbf{Presentation:}

\begin{itemize}
\tightlist
\item
  \textbf{Duration:}~Each group has to submit a 4-5 minutes recording of
  their presentation and demo and upload it into their GitHub
  repository. Please practice before recording to ensure you stay within
  this time frame.
\item
  \textbf{Content:}~Your presentation should include the following
  elements:

  \begin{itemize}
  \tightlist
  \item
    \textbf{Brief Introduction:}~Everyone should have the camera on
    during the recording. A simple way of doing that is to set up a ZOOM
    session with the group. A quick overview of your project, including
    the key features and objectives of your system.
  \item
    \textbf{Demo:}~A live demonstration of your system in action. Focus
    on the core functionalities that highlight your
    system\textquotesingle s uniqueness and adherence to the~project
    requirements.
  \item
    \textbf{Key Challenges \& Learning:}~Briefly discuss any significant
    challenges you faced and what you learned from overcoming them.
  \end{itemize}
\item
\end{itemize}

\paragraph{\texorpdfstring{\textbf{Technical
Setup:}}{Technical Setup:}}\label{technical-setup}

\begin{itemize}
\tightlist
\item
  \textbf{System Preparation:}~Ensure that your application is fully set
  up on the device you will be using for the presentation. This includes
  installing all necessary packages and dependencies required to run
  your application smoothly.
\item
  \textbf{Data Population:}~Populate your application with enough data
  to effectively demonstrate its features. This should include various
  examples that showcase the system\textquotesingle s functionalities
  and capabilities.
\item
  \textbf{Test Beforehand:}~Conduct a thorough test of your system and
  any multimedia components before your presentation to avoid technical
  glitches.
\end{itemize}

\begin{enumerate}
\def\labelenumi{\arabic{enumi}.}
\tightlist
\item
\end{enumerate}

\paragraph{\texorpdfstring{\textbf{Evaluation
Criteria:}}{Evaluation Criteria:}}\label{evaluation-criteria}

Your presentation will be evaluated based on the following:

\begin{itemize}
\tightlist
\item
  \textbf{Adherence to Time Limit}
\item
  \textbf{Clarity and Demonstration of Key Features}
\item
  \textbf{Problem-Solving and Creativity}
\item
  \textbf{Overall Impact}
\end{itemize}

\textbf{Grading Criteria:}

\begin{itemize}
\tightlist
\item
  \textbf{The endpoints must all work properly.}
\item
  \textbf{Unit tests must pass.}
\item
  \textbf{Quality of sprint backlog and tasks.}
\item
  \textbf{Sprint Review must be filled out correctly.}
\item
  \textbf{Attendance (participation in group meetings).}
\end{itemize}

**~One submission per group is enough.**

Points

70

Submitting

a website url or a file upload

\begin{longtable}[]{@{}llll@{}}
\toprule\noalign{}
DueForAvailable fromUntil & & & \\
\midrule\noalign{}
\endhead
\bottomrule\noalign{}
\endlastfoot
Dec 6Dec 6 at 11:59pm & Everyone & -N/A & Dec 7 at 11:59pmDec 7 at
11:59pm \\
\end{longtable}

\begin{center}\rule{0.5\linewidth}{0.5pt}\end{center}

\section{Group Project -
Documentation}\label{group-project---documentation}

\textbf{Your task is to create two distinct documents for the Online
Restaurant Ordering System API: a User Manual and a Technical Document.
These documents are critical for ensuring usability and maintainability
of your API.}

\textbf{User Manual:}

\textbf{Purpose:}~The User Manual is designed to help end-users
understand how to interact with your API effectively.

\textbf{Content Requirements:}

\begin{itemize}
\tightlist
\item
  \textbf{Overview:}~A brief introduction to the API, including its
  purpose and primary features.
\item
  \textbf{Setup Instructions:}~Step-by-step guidance on how to set up
  and start using the API.
\item
  \textbf{Usage Examples:}~Clear examples showing how to make requests
  to the API (including expected input formats) and what responses to
  expect.
\end{itemize}

The User Manual should be user-friendly, easy to navigate, and written
in clear, non-technical language.

\paragraph{\texorpdfstring{\textbf{Technical
Document:}}{Technical Document:}}\label{technical-document}

\textbf{Purpose:}~The Technical Document is intended for future
developers who will maintain, update, or integrate with your API.

\textbf{Content Requirements:}

\begin{itemize}
\tightlist
\item
  \textbf{Architecture Overview:}~A high-level description of the API's
  architecture (include what you did in Part2).
\item
  \textbf{Endpoint Documentation:}~Detailed information about each API
  endpoint, including request methods, parameters, sample requests, and
  expected responses
  (\href{https://instructure.charlotte.edu/courses/264138/files/33808130?wrap=1}{this
  document}\href{https://instructure.charlotte.edu/courses/264138/files/33808130/download?download_frd=1}{Download
  this document}~from Assignment 5 may provide you with some ideas).
\item
  \textbf{Code Examples:}~Snippets or links to significant pieces of
  code, with explanations of their functionality.
\item
  \textbf{Development Environment:}~Instructions for setting up a
  development environment, including required tools, frameworks, and
  dependencies.
\end{itemize}

\textbf{Format:}~The Technical Document should be detailed and precise,
using technical language appropriate for experienced developers. You can
use diagrams, code blocks, and tables to clarify complex points.

\paragraph{\texorpdfstring{\textbf{Submission
Instructions:}}{Submission Instructions:}}\label{submission-instructions}

\begin{itemize}
\tightlist
\item
  Submit both documents in a pdf format.
\item
  Ensure that the documents are well-organized, free of grammatical
  errors, and easy to follow.
\item
  Include a cover page for each document with the API name, your team
  name, and date.
\end{itemize}

Points

15

Submitting

a file upload

\begin{longtable}[]{@{}llll@{}}
\toprule\noalign{}
DueForAvailable fromUntil & & & \\
\midrule\noalign{}
\endhead
\bottomrule\noalign{}
\endlastfoot
Dec 6Dec 6 at 11:59pm & Everyone & -N/A & - \\
\end{longtable}

\end{document}
